\section{Future Directions and Conclusion }
\label{sec:conclusion}

\textbf{Future directions.}
Despite the progress made with \lang, several limitations remain that reveal promising directions for future research. These include extending \lang to support additional output modalities, adapting RadixAttention to operate across multiple levels of the memory hierarchy (e.g., DRAM, Disk)~\cite{sheng2023flexgen}, enabling fuzzy semantic matching within RadixAttention, providing higher-level primitives atop \lang,
fixing starvation in cache-aware scheduling~\cite{sheng2023fairness},
and enhancing the \lang compiler to perform advanced static optimizations such as scheduling and memory planning.

\textbf{Conclusion.}
We introduce \lang, a framework for efficient programming and executing structured language model programs. \lang significantly improves the throughput and latency of complex LM programs through novel optimizations like RadixAttention, compressed finite state machines, and a language interpreter. It is a valuable tool for developing advanced prompting techniques and agent workflows. The source code is publicly available.


\section*{Acknowledgement}
This project is supported by the Stanford Center for Automated Reasoning and gifts from Astronomer, Google, IBM, Intel, Lacework, Microsoft, Mohamed Bin Zayed University of Artificial Intelligence, Nexla, Samsung SDS, Uber, and VMware.
Lianmin Zheng is supported by a Meta Ph.D. Fellowship.
We thank Yuanhan Zhang and Bo Li for the LLaVA-NeXT (video) support.